\section*{A Boolean granularity refinement in exact Fourier degree}

All logarithms are base two.  Let
\(f:\{-1,1\}^n\to\{-1,1\}\) have Fourier degree at most \(d\geq1\),
write \(p_S=\widehat f(S)^2\), and put
\[
 v=\sum_{S\ne\varnothing}p_S=1-p_\varnothing.
\]

\paragraph{Theorem.}
For every such Boolean function,
\[
 \boxed{\quad
 H(f)\leq 2(d-1)v+p_\varnothing\log_2\frac1{p_\varnothing}
 \leq \bigl(2(d-1)+\log_2 e\bigr)I(f).
 \quad}
\]
The first inequality is sharp for the degree-\(d\) address function.

\paragraph{Proof.}
The Boolean granularity lemma says
\[
 \widehat f(S)\in 2^{1-d}\mathbb Z.
\]
For completeness, set
\[
 g(z)=\frac{1-f(1-2z_1,\ldots,1-2z_n)}2.
\]
The multilinear expansion
\(g(z)=\sum_{|T|\leq d}c_T\prod_{i\in T}z_i\) has integer
coefficients by Boolean M\"obius inversion.  Substituting
\(z_i=(1-x_i)/2\) proves the displayed divisibility.
Consequently every nonzero spectral atom satisfies
\(p_S\geq2^{2-2d}\), and hence, for \(S\ne\varnothing\),
\[
 p_S\log_2\frac1{p_S}\leq 2(d-1)p_S.
\]
Summing and retaining the empty atom proves the first inequality.  Next,
\[
 p_\varnothing\log_2\frac1{p_\varnothing}
 =(1-v)\log_2\frac1{1-v}\leq v\log_2e,
\]
while \(v\leq I(f)\), because each nonempty Fourier set has size at
least one.  This proves the second inequality.  For the address function
there are exactly \(2^{2(d-1)}\) nonzero coefficients, all of magnitude
\(2^{1-d}\), and \(p_\varnothing=0\); equality holds in the first bound.
\qed

\paragraph{Analytic location of the loss.}
For arbitrary real degree-\(d\) functions, the Bohnenblust--Hille route
at exponent \(2d/(d+1)\) carries a factor
\(C^{\sqrt{d\log d}}\).  After conversion to R\'enyi entropy of order
\(d/(d+1)\), that factor produces an additional
\(O(d^{3/2}\sqrt{\log d})\) entropy allowance.  Boolean granularity
replaces it by the exact dimension-free support law
\(|\operatorname{supp}\widehat f|\leq2^{2(d-1)}\), or entropy
\(2(d-1)\).  Separating the empty atom then makes the allowance
proportional to \(v\leq I(f)\).

This is a degree-dependent restricted FEI theorem, not full FEI.  The
constant still grows linearly with \(d\).
